\bookchapter{Глава 10: Мобильные платежи и бытовые услуги}

\chapterintro{Честно говоря, поколение моих родителей поначалу сопротивлялось идее «платить телефоном». «Деньги в телефоне — а если украдут?» Это нормальное беспокойство, но реальность такова: наличные в кармане потерять гораздо вероятнее, чем стать жертвой мошенничества с мобильным платежом. Мобильный платёж — это не хранение денег в телефоне, это всего лишь ключ, ключ к вашему банковскому счёту.}

\booksection{Как привязать банковскую карту?}

Неважно, используете ли вы Alipay или WeChat — при первой привязке карты нужно пройти следующие шаги:

\begin{enumerate}
  \item Откройте приложение, найдите «Моё» или «Кошелёк». Нажмите «Добавить банковскую карту».
  \item Введите номер карты. Сейчас большинство приложений поддерживают фото-распознавание номера карты — не нужно вбивать цифры вручную.
  \item Банк отправит SMS с проверочным кодом на ваш зарегистрированный номер телефона. Введите код и установите 6-значный платёжный пароль.
  \item Этот платёжный пароль, ваш банковский PIN-код и пароль блокировки экрана — лучше всего, чтобы это были три разных пароля; не ленитесь и не делайте их одинаковыми.
\end{enumerate}

Маленькая хитрость: если беспокоитесь о безопасности, откройте в банковском приложении электронный счёт класса II или III специально для привязки к мобильным платежам. Держите на этой карте только повседневные карманные деньги — даже если что-то случится, потери будут ограничены.

\booksection{Код оплаты: «кто видит — тот и получает деньги»}

Многие путают код оплаты и код получения денег, а разница между ними очень важна.

\begin{itemize}
  \item \textbf{Код оплаты}: тот штрихкод + QR-код, который вы показываете кассиру. Этот код — ваш «электронный кошелёк»: кто его отсканирует, тот и спишет деньги с вашего счёта. Поэтому\textbf{никогда не делайте скриншот и не отправляйте его никому}. Код автоматически обновляется каждую минуту именно для предотвращения кражи.
  \item \textbf{Код получения денег}: другие люди сканируют этот код, чтобы перевести вам деньги. Этот код можно сохранять как скриншот и отправлять другим — это безопасно.
\end{itemize}

Запомните это различие — оно поможет избежать множества мошеннических ловушек. Любой, кто просит вас «сделать скриншот кода оплаты и прислать» — стопроцентный мошенник.

\booksection{Эти настройки безопасности за 5 минут сэкономят вам массу нервов}

\begin{itemize}
  \item \textbf{Оплата по отпечатку/лицу}: в настройках платежей Alipay и WeChat включите биометрическую оплату. Тогда каждый платёж потребует отпечаток пальца или сканирование лица — даже если телефон потеряется, нашедший не сможет потратить ваши деньги.
  \item \textbf{Замок безопасности}: и в Alipay, и в WeChat есть функция «Замок безопасности». После включения для входа на страницу «Кошелёк» приложения потребуется дополнительная проверка личности. Путь: Alipay Моё Настройки Аккаунт и безопасность Центр безопасности Замок безопасности.
  \item \textbf{Проверка платежей без подтверждения}: многие подписались на разные автоплатежи и забыли об этом. Зайдите в «Настройки платежей» Alipay и WeChat, найдите «Платежи без подтверждения / Автосписание», проверьте, на какие услуги вы подписаны, и отключите ненужные.
  \item \textbf{Управление устройствами входа}: регулярно проверяйте, какие устройства входили в ваш аккаунт. Если обнаружите незнакомое устройство — немедленно отключите. Путь: Alipay Настройки Аккаунт и безопасность Управление устройствами входа.
\end{itemize}

\booksection{Автобус и метро: телефон как транспортная карта}

Сейчас мало кто пользуется физической транспортной картой. Для прохода через турникет телефоном есть два пути:

\begin{itemize}
  \item \textbf{Виртуальная транспортная карта в кошельке телефона}: откройте виртуальную транспортную карту прямо в Apple Wallet, Huawei Wallet или Xiaomi Wallet. Преимущество: не нужно разблокировать телефон — просто приложите, скорость отклика максимальная, работает даже без интернета. В некоторых городах виртуальные карты дают такие же скидки на пересадки, как и физические.
  \item \textbf{QR-код для проезда в Alipay/WeChat}: найдите в приложении «Код для проезда», активируйте для своего города. Преимущество: не нужно пополнять — деньги списываются прямо с баланса или привязанной карты. Но требуется интернет, и перед проходом нужно открыть приложение и найти код.
\end{itemize}

Лично я рекомендую виртуальную транспортную карту в кошельке телефона. Хотя при первом использовании нужно положить 10-20 юаней, пользоваться гораздо удобнее — особенно когда спешишь на метро и не нужно останавливаться, чтобы разблокировать телефон и искать приложение.

\booksection{Коммунальные платежи, доставка еды, такси}

\begin{itemize}
  \item \textbf{Оплата воды, электричества, газа}: в разделе «Бытовые платежи» Alipay или WeChat привяжите номера лицевых счетов воды, электричества и газа. Потом каждый месяц при получении счёта просто заходите и оплачивайте. Номер счёта обычно можно найти на бумажной квитанции или узнать по телефону службы поддержки.
  \item \textbf{Сравнение цен на доставку еды}: цены и скидки на одно и то же блюдо в Meituan и Ele.me часто различаются. Перед заказом загляните в оба приложения — иногда можно сэкономить несколько юаней. Кстати, функция «Совместный заказ» в Meituan позволяет заказывать с друзьями и заметно снижает цену.
  \item \textbf{Безопасность такси}: в Didi или Gaode Maps настройте «Экстренный контакт» и «Поделиться поездкой». При заказе такси приложение автоматически отправит ваш маршрут указанному контакту. Кроме того, поздно ночью лучше выбирать «Такси», а не «Экспресс» — такси находится под отраслевым надзором, что безопаснее.
\end{itemize}

\booksection{Медицинская страховка, соцстрахование и госуслуги через телефон}

За последние годы многое изменилось: то, что раньше требовало походов в учреждения и очередей, теперь делается на телефоне за несколько минут.

\begin{itemize}
  \item \textbf{Электронный полис медицинского страхования}: найдите в Alipay или WeChat «Электронный полис медстрахования», активируйте — и при посещении врача или покупке лекарств в аптеке просто покажите телефон, физическая карта не нужна. При первой активации требуется верификация по лицу, затем при каждом использовании генерируется динамический код.
  \item \textbf{Проверка соцстрахования}: скачайте приложение «Zhangshang 12333» или используйте «Городской центр» в Alipay — можно проверить историю отчислений, баланс пенсионного счёта, подать заявление на пособие по безработице и т.д.
  \item \textbf{Подоходный налог}: ежегодная сверка НДФЛ с марта по июнь делается в приложении «Налог на доходы физических лиц». Проверьте, есть ли у вас право на специальные дополнительные вычеты (образование детей, проценты по ипотеке, содержание пожилых родителей и т.д.) — они помогут уменьшить налог или получить возврат.
\end{itemize}

\begin{tcbwarn}{Запомните эти правила — защитите себя от мошенников}
\begin{itemize}
  \item Правоохранительные органы, банки и службы поддержки никогда не попросят вас перевести деньги на «безопасный счёт».
  \item Любой, кто просит вас сделать скриншот кода оплаты и прислать — мошенник.
  \item Не переходите по незнакомым ссылкам, особенно в SMS про «просроченный ETC» или «заблокированную карту соцстрахования».
  \item Возврат средств за онлайн-покупки происходит автоматически по тому же пути — вам не нужно никому сообщать номер банковской карты.
\end{itemize}

\end{tcbwarn}

\begin{tcbpractice}{Сделайте прямо сейчас: проверка безопасности ваших платёжных приложений}
\begin{enumerate}
  \item Откройте Alipay и WeChat, убедитесь, что включена оплата по отпечатку/лицу.
  \item Найдите функцию «Замок безопасности» и включите её.
  \item Проверьте список «Платежи без подтверждения / Автосписание» и отключите ненужные услуги.
  \item Проверьте список «Устройства входа» и удалите незнакомые устройства.
  \item Активируйте электронный полис медстрахования — в следующий раз в больницу без физической карты.
\end{enumerate}

\end{tcbpractice}
